%%%%%%%%%%%%%%%%%%%%%%%%%%%%%%%%%%%%%%%%%%%%%%%%%%%%%%%%%%%%%%%%%%%%%%%%
% Piezas de dibujo de los péndulos acoplados (TikZ)
%
% Vista frontal y minimalista del montaje de la práctica: dos cabezales
% triangulares (como los soportes PHYWE de las fotos), dos varillas de
% longitud L con su masa lenticular y un muelle helicoidal sujeto a las dos
% varillas a una distancia \ell de los ejes de giro.
% Se usan en la portada y en la figura del montaje experimental.
% Requiere los colores azulprof, azulacero, turquesa, azulclaro,
% grispizarra y grisclaro (definidos en main.tex).
%%%%%%%%%%%%%%%%%%%%%%%%%%%%%%%%%%%%%%%%%%%%%%%%%%%%%%%%%%%%%%%%%%%%%%%%

% Punto a distancia #2 del pivote #1 a lo largo de una varilla girada #3 grados
\newcommand{\sobrevarilla}[3]{($(#1)+({#2*sin(#3)},{-#2*cos(#3)})$)}

% Cabezal triangular invertido con el eje de giro en el vértice inferior
% \cabezal{coordenada}{escala}
\newcommand{\cabezal}[2]{%
  \fill[azulprof] ($(#1)+(-0.62*#2,0.78*#2)$) -- ($(#1)+(0.62*#2,0.78*#2)$)
    -- ($(#1)+(0,-0.06*#2)$) -- cycle;
  \fill[azulacero!70!azulprof] ($(#1)+(-0.62*#2,0.78*#2)$) -- ($(#1)+(0.62*#2,0.78*#2)$)
    -- ($(#1)+(0.62*#2,0.86*#2)$) -- ($(#1)+(-0.62*#2,0.86*#2)$) -- cycle;
  \draw[white, line width=0.5pt, fill=azulprof] (#1) circle (0.085*#2);
  \fill[white] ($(#1)+(0,0.46*#2)$) circle (0.05*#2);}

% Muelle helicoidal entre dos puntos, con sus pinzas
% \muelleacople{A}{B}{escala}
\newcommand{\muelleacople}[3]{%
  \draw[azulacero, line width=0.55pt,
        decorate, decoration={coil, aspect=0.55, amplitude=3.1pt*#3,
                              segment length=2.9pt*#3, pre length=3pt*#3, post length=3pt*#3}]
    (#1) -- (#2);
  \fill[grispizarra] ($(#1)+(-0.07*#3,-0.07*#3)$) rectangle ($(#1)+(0.07*#3,0.07*#3)$);
  \fill[grispizarra] ($(#2)+(-0.07*#3,-0.07*#3)$) rectangle ($(#2)+(0.07*#3,0.07*#3)$);}

% Masa lenticular: disco con un brillo discreto
% \masalenticular{coordenada}{radio}
\newcommand{\masalenticular}[2]{%
  \shade[ball color=turquesa!85!white, opacity=1] (#1) circle (#2);
  \draw[azulprof, line width=0.6pt] (#1) circle (#2);
  \fill[white, opacity=0.35] ($(#1)+(-0.32*#2,0.34*#2)$) circle (0.22*#2);}

% El par de péndulos acoplados
% \parpendulos{phi1}{phi2}{ell}{L}{d}{escala}
%   phi1, phi2 : ángulos (grados) respecto a la vertical, positivos hacia la derecha
%   ell, L     : distancia del eje al muelle y longitud de la varilla
%   d          : separación entre los ejes de giro
%   escala     : tamaño de cabezales, masas y muelle
% Define las coordenadas P1, P2 (ejes), S1, S2 (muelle) y B1, B2 (masas).
\newcommand{\parpendulos}[6]{%
  \coordinate (P1) at ({-#5/2},0);
  \coordinate (P2) at ({#5/2},0);
  \coordinate (S1) at ($(P1)+({#3*sin(#1)},{-#3*cos(#1)})$);
  \coordinate (S2) at ($(P2)+({#3*sin(#2)},{-#3*cos(#2)})$);
  \coordinate (B1) at ($(P1)+({#4*sin(#1)},{-#4*cos(#1)})$);
  \coordinate (B2) at ($(P2)+({#4*sin(#2)},{-#4*cos(#2)})$);
  % verticales de equilibrio
  \draw[grispizarra!45, line width=0.45pt, dash pattern=on 3pt off 2.2pt]
    (P1) -- ++(0,{-#4-0.55*#6});
  \draw[grispizarra!45, line width=0.45pt, dash pattern=on 3pt off 2.2pt]
    (P2) -- ++(0,{-#4-0.55*#6});
  % varillas
  \draw[azulprof, line width=1.05pt*#6, line cap=round] (P1) -- (B1);
  \draw[azulprof, line width=1.05pt*#6, line cap=round] (P2) -- (B2);
  % muelle de acoplamiento
  \muelleacople{S1}{S2}{#6}
  % masas lenticulares y cabezales
  \masalenticular{B1}{0.36*#6}
  \masalenticular{B2}{0.36*#6}
  \cabezal{P1}{#6}
  \cabezal{P2}{#6}}

% Arco del ángulo girado por una varilla, con su rótulo. El rótulo se coloca
% al otro lado de la vertical de equilibrio, para que no pise la varilla.
% \angulo{pivote}{phi}{radio}{rótulo}
\newcommand{\angulo}[4]{%
  \draw[turquesa!80!azulprof, line width=0.6pt, -{Stealth[length=3.4pt]}]
    ($(#1)+(0,-#3)$) arc[start angle=-90, end angle={-90+#2}, radius=#3];
  \ifdim #2pt>0pt
    \node[text=azulprof, font=\small, anchor=east, inner sep=1.5pt] at ($(#1)+(-0.06,{-#3})$) {#4};
  \else
    \node[text=azulprof, font=\small, anchor=west, inner sep=1.5pt] at ($(#1)+(0.06,{-#3})$) {#4};
  \fi}

% Posición «fantasma» de un péndulo (el extremo opuesto de su oscilación)
% \fantasma{pivote}{phi}{L}{escala}
\newcommand{\fantasma}[4]{%
  \draw[azulclaro!90!grispizarra, line width=0.8pt*#4, line cap=round, opacity=0.8]
    (#1) -- ($(#1)+({#3*sin(#2)},{-#3*cos(#2)})$);
  \draw[azulclaro!90!grispizarra, line width=0.5pt, fill=azulfondo, opacity=0.9]
    ($(#1)+({#3*sin(#2)},{-#3*cos(#2)})$) circle (0.36*#4);}

% Flecha de movimiento a lo largo de la trayectoria de la masa
% \movimiento{pivote}{desde}{hasta}{radio}  (ángulos respecto a la vertical)
\newcommand{\movimiento}[4]{%
  \draw[turquesa, line width=0.75pt, -{Stealth[length=3.6pt]}]
    ($(#1)+({-90+#2}:#4)$) arc[start angle={-90+#2}, end angle={-90+#3}, radius=#4];}

% Pequeña pulsación: viñeta de la paginación y ornamento entre secciones
\newcommand{\ondita}[1][1]{%
  \tikz[baseline=-0.55ex, x=#1*1pt, y=#1*1pt]{%
    \draw[turquesa, line width=0.45pt, line join=round]
      plot[domain=0:18, samples=90, smooth, variable=\t]
      ({\t},{2.6*cos(2.1*\t r)*sin(\t*180/18)});
    \draw[azulprof!60, line width=0.3pt, densely dotted]
      plot[domain=0:18, samples=30, smooth, variable=\t] ({\t},{2.6*sin(\t*180/18)});}}
